% AlgoTrace: research paper.
% Build: latexmk in the paper/ folder (settings in .latexmkrc); works with pdfLaTeX and LuaLaTeX, bibliography: biber.
\documentclass[a4paper,11pt]{article}

% ---------- Fonts and encoding (pdfLaTeX or LuaLaTeX) ----------
\usepackage{iftex}
\ifPDFTeX
  \usepackage[T1]{fontenc}
  \usepackage[utf8]{inputenc}
  \usepackage{lmodern}
\else
  \usepackage{fontspec}
\fi

% ---------- Language ----------
% Main language of the text: main=english (team decision); the abstract is in both languages.
\usepackage[polish,main=english]{babel}
\usepackage{csquotes}

% ---------- Page layout and typography ----------
\usepackage[margin=2.5cm]{geometry}
\usepackage{microtype}
\usepackage{xcolor}

% ---------- Mathematics, tables, figures, algorithms ----------
\usepackage{amsmath}
\usepackage{amssymb}
\usepackage{booktabs}
\usepackage{graphicx}
\graphicspath{{figures/}}
\usepackage{tikz}
\usetikzlibrary{positioning,arrows.meta,babel}
\usepackage{listings}
\lstset{basicstyle=\ttfamily\small,columns=fullflexible,breaklines=true,frame=single}

% ---------- Bibliography ----------
\usepackage[backend=biber,style=numeric,sorting=none,maxbibnames=10]{biblatex}
\addbibresource{references.bib}

% ---------- Hyperlinks (near the end of the preamble) ----------
\usepackage[hidelinks]{hyperref}

% ---------- Pseudocode (after hyperref) ----------
\usepackage{algorithm}
\usepackage{algpseudocode}
% Polish name of the algorithm float, used only inside otherlanguage{polish}.
\addto\captionspolish{\floatname{algorithm}{Algorytm}}

% ---------- Project commands ----------
\newcommand{\TODO}[1]{\textcolor{red}{[TODO: #1]}}
\newcommand{\tool}[1]{\texttt{#1}}

% ---------- Metadata ----------
\title{AlgoTrace: Which Algorithms Do Language Models Use to Solve Tasks Through Tool Calling?}
\author{AlgoTrace project team\\
  Koło Naukowe Informatyków, University of Rzeszów\\
  \TODO{authors and affiliations}}
\date{\today}

\begin{document}

\maketitle

% Streszczenia w obu językach, niezależnie od języka głównego (main= w main.tex).
\begin{otherlanguage}{polish}
\begin{abstract}
  \TODO{streszczenie po polsku, ok. 150--250 słów: problem, metoda, wyniki, wnioski.}
\end{abstract}
\noindent\textbf{Słowa kluczowe:} modele językowe, tool calling, algorytmy sortowania, analiza przebiegów.
\end{otherlanguage}

\begin{otherlanguage}{english}
\begin{abstract}
  \TODO{abstract in English, about 150--250 words.}
\end{abstract}
\noindent\textbf{Keywords:} large language models, tool calling, sorting algorithms, trace analysis.
\end{otherlanguage}


\tableofcontents

\section{Wstęp}
\label{sec:introduction}

\TODO{problem: modele językowe działają jako agenci przez wywołania narzędzi, a testy tool callingu oceniają głównie
poprawność wywołań i wyniku, rzadko to, jak model doszedł do rozwiązania.}

\subsection{Cel badania}
\label{sec:goal}

\TODO{cel: otwarte środowisko testowe z analizatorem przebiegów i raport, jak skutecznie i jakimi algorytmami modele
LLM różnych dostawców rozwiązują zadania przez tool calling.}

\subsection{Pytania badawcze}
\label{sec:questions}

\begin{enumerate}
  \item \textbf{Skuteczność:} jak odsetek poprawnie rozwiązanych zadań zależy od modelu, dostawcy i rozmiaru
    zadania $N$?
  \item \textbf{Algorytm:} jakie algorytmy modele faktycznie wykonują: znane, ich mieszanki czy własne strategie?
  \item \textbf{Deklaracja a wykonanie:} czy algorytm zapowiedziany w tekście zgadza się z wykonanym?
  \item \textbf{Efektywność:} ile porównań i zamian zużywa model wobec granic teoretycznych i algorytmów wzorcowych?
  \item \textbf{Zmiana strategii:} czy model zmienia algorytm, gdy rośnie $N$ albo gdy narzędzie zwróci błąd?
\end{enumerate}

\subsection{Wkład}
\label{sec:contribution}

\TODO{co nowego: każdy ruch jako osobne wywołanie narzędzia na ukrytym stanie, rozpoznawanie algorytmu z zapisu
ruchów, porównanie deklaracji z wykonaniem, gotowe modele wielu dostawców bez trenowania.}

\section{Powiązane prace}
\label{sec:related-work}

\TODO{rozwinąć po przeglądzie literatury; punkty wyjścia poniżej.}

\begin{itemize}
  \item Łamigłówki, w tym wieże Hanoi, jako test modeli z rozumowaniem; model podaje całą listę ruchów w jednej
    odpowiedzi, bez narzędzi~\cite{shojaee2025illusion}.
  \item Polemika: limit długości odpowiedzi jako część przyczyny załamania wyników~\cite{lawsen2025comment}.
  \item Te same łamigłówki z narzędziami (interpreter, notatnik)~\cite{song2025thinking}.
  \item Model uczony ze wzmocnieniem sortuje akcjami porównania i zamiany na częściowo ukrytym stanie; najbliższa
    praca do naszej~\cite{yue2025earl}.
  \item \TODO{benchmarki tool callingu i prace o rozpoznawaniu algorytmu z ciągu operacji.}
\end{itemize}

\section{Metoda}
\label{sec:method}

\subsection{Środowisko}
\label{sec:environment}

Model nie widzi danych zadania: zmienia i odczytuje ukryty stan tylko przez narzędzia, a każde wywołanie trafia do
zapisu przebiegu (rys.~\ref{fig:environment}).

\begin{figure}[htbp]
  \centering
  \begin{tikzpicture}[
      node distance=6mm and 7mm,
      box/.style={draw, rounded corners, align=center, font=\small, minimum height=10mm, minimum width=24mm},
      arr/.style={-{Stealth}, thick},
      both/.style={{Stealth}-{Stealth}, thick}]
    \node[box] (gen) {Generator\\zadań};
    \node[box, right=of gen] (sandbox) {Piaskownica\\(ukryty stan)};
    \node[box, right=of sandbox] (loop) {Pętla\\tool callingu};
    \node[box, right=of loop] (model) {Model LLM\\(przez adapter)};
    \node[box, below=of loop] (trace) {Zapis\\przebiegu};
    \node[box, left=of trace] (analyzer) {Analizator};
    \node[box, left=of analyzer] (refs) {Algorytmy\\wzorcowe};
    \draw[arr] (gen) -- (sandbox);
    \draw[both] (sandbox) -- (loop);
    \draw[both] (loop) -- (model);
    \draw[arr] (loop) -- (trace);
    \draw[arr] (trace) -- (analyzer);
    \draw[arr] (refs) -- (analyzer);
  \end{tikzpicture}
  \caption{Przepływ danych: wykonanie zadania (góra) i analiza po fakcie (dół).}
  \label{fig:environment}
\end{figure}

\TODO{generator z ziarnem, walidacja argumentów, limity, wstrzykiwanie błędów, adaptery dostawców.}

\subsection{Zadania}
\label{sec:tasks}

\begin{table}[htbp]
  \centering
  \caption{Rodziny zadań, narzędzia i granice teoretyczne.}
  \label{tab:tasks}
  \small
  \begin{tabular}{@{}p{2.8cm}p{3.4cm}p{4.2cm}p{3.6cm}@{}}
    \toprule
    Zadanie & Narzędzia & Wzorce & Granica \\
    \midrule
    Sortowanie $N$ elementów & \tool{compare(i, j)}, \tool{swap(i, j)} & bąbelkowe, przez wstawianie, przez wybór,
      szybkie, scalanie & $\lceil \log_2 N! \rceil$ porównań \\
    Wieże Hanoi, $n$ krążków & \tool{move(from, to)} & rekurencyjne & $2^n - 1$ ruchów \\
    Wyszukiwanie & \tool{read(index)} & binarne, liniowe, interpolacyjne & $\lceil \log_2 (N + 1) \rceil$ odczytów
      (binarne) \\
    \bottomrule
  \end{tabular}
\end{table}

\subsection{Rozpoznawanie algorytmu}
\label{sec:recognition}

\TODO{cechy przebiegu, odtworzenie wzorców na tych samych danych, odległość ciągów operacji, klasyfikacja: znany
algorytm, hybryda, strategia nowa; walidacja na syntetycznych przebiegach.}

\begin{algorithm}[htbp]
  \caption{Klasyfikacja przebiegu sortowania (szkic).}
  \label{alg:classify}
  \begin{algorithmic}[1]
    \Require zapis przebiegu $T$, dane wejściowe $x$, warianty wzorców $\mathcal{W}$, próg $\tau$
    \If{stan końcowy nie jest posortowany}
      \State \Return błąd i jego rodzaj
    \EndIf
    \ForAll{$w \in \mathcal{W}$}
      \State $d_w \gets$ odległość ciągu operacji $T$ od przebiegu $w(x)$
    \EndFor
    \If{$\min_w d_w < \tau$}
      \State \Return znany algorytm $\arg\min_w d_w$
    \ElsIf{fragmenty $T$ pasują do różnych wzorców}
      \State \Return hybryda
    \Else
      \State \Return strategia nowa
    \EndIf
  \end{algorithmic}
\end{algorithm}

\subsection{Miary}
\label{sec:measures}

\TODO{skuteczność; porównania wobec $\lceil \log_2 N! \rceil$ i wzorca, porównania zbędne na grafie relacji
elementów; zamiany wobec liczby inwersji i $N - c$, gdzie $c$ to liczba cykli permutacji; błędy; zgodność
deklaracji; tokeny, czas i koszt.}

\section{Experimental protocol}
\label{sec:experiment}

\subsection{Models and providers}
\label{sec:models}

\TODO{table: provider, model, version identifier, date of the runs, local or cloud model.}

\subsection{Runs and settings}
\label{sec:settings}

\TODO{prompt and task description, algorithm declaration, model settings, limits, task sizes, repetitions, seeds;
following the protocol in docs/protocol.}

\subsection{Analysis plan}
\label{sec:analysis-plan}

\TODO{confidence intervals, comparison tests, classifier threshold and hybrid rule.}

\section{Results}
\label{sec:results}

% Tables and plots are generated into figures/ by the analysis scripts; they are not drawn by hand.

\subsection{Analyzer validation}
\label{sec:validation}

\TODO{accuracy on synthetic traces, confusion matrix.}

\subsection{Success}
\label{sec:success}

\TODO{research question 1.}

\subsection{Recognized algorithms}
\label{sec:algorithms}

\TODO{research questions 2 and 3.}

\subsection{Efficiency and change of strategy}
\label{sec:efficiency}

\TODO{research questions 4 and 5.}

\section{Discussion}
\label{sec:discussion}

\TODO{interpretation of the results against related work.}

\subsection{Limitations and threats to validity}
\label{sec:threats}

\begin{itemize}
  \item \TODO{the prompt may suggest an algorithm.}
  \item \TODO{with a small $N$, different algorithms leave similar traces.}
  \item \TODO{non-determinism and silent model updates.}
  \item \TODO{the conclusions apply only to the models and tasks studied.}
\end{itemize}

\section{Podsumowanie}
\label{sec:conclusion}

\TODO{odpowiedzi na pytania badawcze w skrócie, udostępnione środowisko i zbiór przebiegów, dalsze prace.}


\printbibliography

\appendix
\section{Kontrakt danych}
\label{app:data-contract}

\TODO{skrót formatów z schemas/: konfiguracja eksperymentu, zapis przebiegu, wiersz wyników, wynik analizatora.}

Przykład zdarzenia zapisu przebiegu (format do ustalenia):

\begin{lstlisting}
{"step": 1, "turn": 1, "tool": "compare", "args": {"i": 0, "j": 1}, "result": -1}
{"step": 2, "turn": 2, "tool": "compare", "args": {"i": 1, "j": 2}, "result": 1}
{"step": 3, "turn": 2, "tool": "swap", "args": {"i": 1, "j": 2}, "result": "ok"}
\end{lstlisting}


\end{document}
